\chapter{El modelo final}\label{cap:modelo-final}

\section{Especificación y coeficientes}

El modelo final, estimado por MCO con \res{n-final} condados de \res{n-estados-final} estados, es
\[
\begin{aligned}
\widehat{\text{Mortalidad}} ={}& \beta_0
 + \beta_1\,\text{Incidencia} + \beta_2\,\text{Edad mediana}
 + \beta_3\,\text{Tamaño del hogar}\\
&+ \beta_4\,\text{Bachillerato 18--24} + \beta_5\,\text{Grado universitario 25+}
 + \beta_6\,\text{Bachillerato 25+}\\
&+ \beta_7\,\text{Desempleo} + \beta_8\,\text{Sólo seguro público}
 + \beta_9\,\text{Otra raza} + \beta_{10}\,\text{Natalidad}\\
&+ \textstyle\sum_{r}\gamma_r\,\text{Región}_r
 + \sum_{r}\delta_r\,(\text{Grado universitario}\times\text{Región}_r)\\
&+ \textstyle\sum_{r}\eta_r\,(\text{Desempleo}\times\text{Región}_r),
\end{aligned}
\]
con las explicativas continuas centradas en su media y $r\in\{\text{Noreste},\text{Medio Oeste},
\text{Oeste}\}$ (el Sur es la referencia): \res{k-final} regresores además de la constante. La
\cref{tab:coeficientes} recoge las estimaciones con errores típicos robustos por conglomerados de
estado e intervalos de confianza con la $t$ de \res{gl-inferencia} grados de libertad.

{\footnotesize
\begin{xltabular}{\linewidth}{>{\raggedright\arraybackslash}X r r r r l r}
\caption{Coeficientes del modelo final de efectos (MCO con errores típicos robustos por conglomerados de estado)}\label{tab:coeficientes} \\
\toprule
Término & $\hat\beta$ & E.\,T. & $t$ & $p$ & IC 95\,\% & Beta \\
\midrule
\endfirsthead
\multicolumn{7}{l}{\emph{(continuación)}} \\
\toprule
Término & $\hat\beta$ & E.\,T. & $t$ & $p$ & IC 95\,\% & Beta \\
\midrule
\endhead
\midrule
\multicolumn{7}{r}{\emph{(continúa)}} \\
\endfoot
\bottomrule
\multicolumn{7}{p{0.97\linewidth}}{\footnotesize $n = \num{2840}$ condados de \num{48} estados. Explicativas continuas centradas en su media: la constante es la mortalidad esperada de un condado medio del Sur. Inferencia con $t$ de \num{47} g.\,l. Beta: coeficiente tipificado.} \\
\endlastfoot
Constante & \num{184.526} & \num{1.319} & \num{139.95} & $<\num{0.001}$ & $[\num{181.873};$\allowbreak\ $\num{187.178}]$ &  \\
Incidencia de cáncer & \num{0.187} & \num{0.019} & \num{9.61} & $<\num{0.001}$ & $[\num{0.148};$\allowbreak\ $\num{0.226}]$ & \num{0.379} \\
Edad mediana & \num{-0.788} & \num{0.187} & \num{-4.22} & $<\num{0.001}$ & $[\num{-1.164};$\allowbreak\ $\num{-0.413}]$ & \num{-0.147} \\
Tamaño medio del hogar & \num{-12.875} & \num{3.292} & \num{-3.91} & $<\num{0.001}$ & $[\num{-19.499};$\allowbreak\ $\num{-6.252}]$ & \num{-0.115} \\
Bachillerato (18–24) & \num{0.241} & \num{0.064} & \num{3.76} & $<\num{0.001}$ & $[\num{0.112};$\allowbreak\ $\num{0.370}]$ & \num{0.078} \\
Grado universitario (25+) & \num{-1.054} & \num{0.255} & \num{-4.14} & $<\num{0.001}$ & $[\num{-1.566};$\allowbreak\ $\num{-0.542}]$ & \num{-0.207} \\
Desempleo (16+) & \num{0.366} & \num{0.284} & \num{1.29} & \num{0.204} & $[\num{-0.205};$\allowbreak\ $\num{0.937}]$ & \num{0.045} \\
Solo seguro público & \num{0.892} & \num{0.214} & \num{4.18} & $<\num{0.001}$ & $[\num{0.463};$\allowbreak\ $\num{1.322}]$ & \num{0.197} \\
Otra raza & \num{-0.660} & \num{0.210} & \num{-3.14} & \num{0.003} & $[\num{-1.082};$\allowbreak\ $\num{-0.238}]$ & \num{-0.084} \\
Natalidad & \num{-0.381} & \num{0.249} & \num{-1.53} & \num{0.132} & $[\num{-0.883};$\allowbreak\ $\num{0.120}]$ & \num{-0.027} \\
Región: Noreste & \num{-15.683} & \num{2.316} & \num{-6.77} & $<\num{0.001}$ & $[\num{-20.343};$\allowbreak\ $\num{-11.023}]$ &  \\
Región: Medio Oeste & \num{-4.327} & \num{2.419} & \num{-1.79} & \num{0.080} & $[\num{-9.194};$\allowbreak\ $\num{0.541}]$ &  \\
Región: Oeste & \num{-14.131} & \num{2.756} & \num{-5.13} & $<\num{0.001}$ & $[\num{-19.675};$\allowbreak\ $\num{-8.586}]$ &  \\
Bachillerato (25+) & \num{0.354} & \num{0.192} & \num{1.85} & \num{0.071} & $[\num{-0.031};$\allowbreak\ $\num{0.739}]$ & \num{0.091} \\
Grado universitario (25+) $\times$ Noreste & \num{0.844} & \num{0.219} & \num{3.86} & $<\num{0.001}$ & $[\num{0.404};$\allowbreak\ $\num{1.285}]$ &  \\
Grado universitario (25+) $\times$ Medio Oeste & \num{0.229} & \num{0.308} & \num{0.74} & \num{0.461} & $[\num{-0.391};$\allowbreak\ $\num{0.849}]$ &  \\
Grado universitario (25+) $\times$ Oeste & \num{-0.031} & \num{0.257} & \num{-0.12} & \num{0.905} & $[\num{-0.547};$\allowbreak\ $\num{0.485}]$ &  \\
Desempleo (16+) $\times$ Noreste & \num{0.903} & \num{0.555} & \num{1.63} & \num{0.110} & $[\num{-0.214};$\allowbreak\ $\num{2.020}]$ &  \\
Desempleo (16+) $\times$ Medio Oeste & \num{1.478} & \num{0.446} & \num{3.31} & \num{0.002} & $[\num{0.581};$\allowbreak\ $\num{2.375}]$ &  \\
Desempleo (16+) $\times$ Oeste & \num{-0.429} & \num{0.688} & \num{-0.62} & \num{0.536} & $[\num{-1.813};$\allowbreak\ $\num{0.955}]$ &  \\
\end{xltabular}
}


\section{Salida tipo SPSS: \texorpdfstring{$R^2$}{R²}, \texorpdfstring{$R^2$}{R²} ajustado y RMSE}\label{sec:salida-spss}

\begin{entregable}{(b): salida del software}
El modelo explica el \res{pct-r2-final}\,\% de la variabilidad de la mortalidad entre condados:
$R^2=\res{r2-final}$ y $\bar R^2=\res{r2aj-final}$. El error típico de la estimación es
$\hat\sigma=\sqrt{\mathrm{SCE}/(n-k-1)}=\res{sigma-final}$ y la raíz del error cuadrático medio
dentro de la muestra, $\mathrm{RMSE}=\sqrt{\mathrm{SCE}/n}=\res{rmse-final}$ muertes por
\num{100000} habitantes. Fuera de la muestra, la validación cruzada de esta especificación da un
RMSE de \res{cv-rmse-mco-efectos} y la partición de prueba, \res{test-rmse-mco-efectos}
(\cref{cap:prediccion}).
\end{entregable}

Las \cref{tab:spss-resumen,tab:anova,tab:spss-coeficientes} reproducen las tablas que imprime el
procedimiento \var{REGRESSION} de SPSS: resumen del modelo, ANOVA y coeficientes con errores
típicos clásicos, coeficientes tipificados (Beta), tolerancia y FIV. La sintaxis del
\cref{ap:reproducibilidad}, generada automáticamente a partir de la corrida, aplica la misma
depuración sobre \var{practica.sav} y obtiene exactamente estas cifras.

\begin{table}[htbp]
\centering
\footnotesize
\caption{Resumen del modelo final (salida tipo SPSS)}\label{tab:spss-resumen}
\begin{tabular}{r r r r r r r r r r}
\toprule
$R$ & $R^2$ & $\bar R^2$ & E.\,T. estimación & RMSE & D-W & $F$ & g.\,l.$_1$ & g.\,l.$_2$ & $p$ \\
\midrule
\num{0.747} & \num{0.558} & \num{0.555} & \num{18.50} & \num{18.44} & \num{1.878} & \num{187.34} & \num{19} & \num{2820} & $<\num{0.001}$ \\
\bottomrule
\end{tabular}
\par\smallskip{\footnotesize\raggedright E.\,T. de la estimación: $\sqrt{\mathrm{SCE}/(n-k-1)}$; RMSE: $\sqrt{\mathrm{SCE}/n}$, ambos dentro de la muestra. D-W: Durbin-Watson. El $F$ y su $p$ son los clásicos que imprime SPSS.\par}
\end{table}

\begin{table}[htbp]
\centering
\small
\caption{Tabla ANOVA del modelo final}\label{tab:anova}
\begin{tabular}{l r r r r r}
\toprule
Fuente & Suma de cuadrados & g.\,l. & Cuadrado medio & $F$ & $p$ \\
\midrule
Regresión (SCR) & \num{1218548.0} & 19 & \num{64134.1} & \num{187.34} & $<\num{0.001}$ \\
Residual (SCE) & \num{965407.2} & \num{2820} & \num{342.34} & & \\
Total (SCT) & \num{2183955.0} & \num{2839} & & & \\
\bottomrule
\end{tabular}
\par\smallskip{\footnotesize\raggedright $R^2 = \num{0.558}$, $\bar R^2 = \num{0.555}$. Contraste global robusto (Wald, cluster): $F = \num{94.3}$, $p < \num{0.001}$.\par}
\end{table}


Recuérdese que
\[
R^2=1-\frac{\mathrm{SCE}}{\mathrm{SCT}},\qquad
\bar R^2=1-\frac{\mathrm{SCE}/(n-k-1)}{\mathrm{SCT}/(n-1)}=1-(1-R^2)\frac{n-1}{n-k-1}.
\]
El $R^2$ ajustado penaliza el número de regresores; con $n$ tan grande frente a $k$ apenas
difiere del $R^2$. Un $R^2$ del orden de 0,56 es alto para datos ecológicos transversales: buena
parte de la variabilidad restante es ruido de muestreo de las tasas de los condados pequeños
(\cref{sec:estructura-error}), que ninguna explicativa puede capturar.

El contraste global clásico ($F=\res{f-clasico}$) y su versión robusta (Wald con covarianza por
conglomerados, $F=\res{f-robusto}$, \res{f-robusto-p}) rechazan con claridad que todos los
coeficientes sean nulos.

{\footnotesize
\begin{xltabular}{\linewidth}{X r r r r r l r r}
\caption{Coeficientes del modelo final con errores típicos clásicos (salida tipo SPSS)}\label{tab:spss-coeficientes} \\
\toprule
Término & $B$ & E.\,T. & Beta & $t$ & Sig. & IC 95\,\% & Tol. & FIV \\
\midrule
\endfirsthead
\multicolumn{9}{l}{\emph{(continuación)}} \\
\toprule
Término & $B$ & E.\,T. & Beta & $t$ & Sig. & IC 95\,\% & Tol. & FIV \\
\midrule
\endhead
\midrule
\multicolumn{9}{r}{\emph{(continúa)}} \\
\endfoot
\bottomrule
\multicolumn{9}{p{0.97\linewidth}}{\footnotesize Errores típicos clásicos $\hat\sigma^2(X^\top X)^{-1}$, como en el procedimiento REGRESSION; la sintaxis \texttt{spss/modelo\_final.sps} los reproduce. La inferencia del informe usa los robustos (tabla~\ref{tab:errores}).} \\
\endlastfoot
Constante & \num{184.526} & \num{0.545} &  & \num{338.37} & $<\num{0.001}$ & $[\num{183.457};\ \num{185.595}]$ &  &  \\
Incidencia de cáncer & \num{0.187} & \num{0.007} & \num{0.379} & \num{27.66} & $<\num{0.001}$ & $[\num{0.173};\ \num{0.200}]$ & \num{0.835} & \num{1.20} \\
Edad mediana & \num{-0.788} & \num{0.095} & \num{-0.147} & \num{-8.31} & $<\num{0.001}$ & $[\num{-0.974};\ \num{-0.602}]$ & \num{0.501} & \num{1.99} \\
Tamaño medio del hogar & \num{-12.875} & \num{2.004} & \num{-0.115} & \num{-6.42} & $<\num{0.001}$ & $[\num{-16.805};\ \num{-8.945}]$ & \num{0.486} & \num{2.06} \\
Bachillerato (18–24) & \num{0.241} & \num{0.047} & \num{0.078} & \num{5.17} & $<\num{0.001}$ & $[\num{0.150};\ \num{0.333}]$ & \num{0.684} & \num{1.46} \\
Grado universitario (25+) & \num{-1.054} & \num{0.160} & \num{-0.207} & \num{-6.60} & $<\num{0.001}$ & $[\num{-1.367};\ \num{-0.741}]$ & \num{0.160} & \num{6.25} \\
Desempleo (16+) & \num{0.366} & \num{0.178} & \num{0.045} & \num{2.06} & \num{0.040} & $[\num{0.017};\ \num{0.715}]$ & \num{0.323} & \num{3.09} \\
Solo seguro público & \num{0.892} & \num{0.093} & \num{0.197} & \num{9.63} & $<\num{0.001}$ & $[\num{0.711};\ \num{1.074}]$ & \num{0.374} & \num{2.67} \\
Otra raza & \num{-0.660} & \num{0.118} & \num{-0.084} & \num{-5.60} & $<\num{0.001}$ & $[\num{-0.891};\ \num{-0.429}]$ & \num{0.702} & \num{1.42} \\
Natalidad & \num{-0.381} & \num{0.187} & \num{-0.027} & \num{-2.04} & \num{0.042} & $[\num{-0.749};\ \num{-0.014}]$ & \num{0.888} & \num{1.13} \\
region[Noreste] & \num{-15.683} & \num{1.685} &  & \num{-9.30} & $<\num{0.001}$ & $[\num{-18.988};\ \num{-12.378}]$ & \num{0.601} & \num{1.66} \\
region[Medio Oeste] & \num{-4.327} & \num{0.957} &  & \num{-4.52} & $<\num{0.001}$ & $[\num{-6.202};\ \num{-2.451}]$ & \num{0.632} & \num{1.58} \\
region[Oeste] & \num{-14.131} & \num{1.221} &  & \num{-11.57} & $<\num{0.001}$ & $[\num{-16.525};\ \num{-11.737}]$ & \num{0.668} & \num{1.50} \\
Bachillerato (25+) & \num{0.354} & \num{0.096} & \num{0.091} & \num{3.68} & $<\num{0.001}$ & $[\num{0.165};\ \num{0.542}]$ & \num{0.257} & \num{3.89} \\
Grado universitario (25+) $\times$ region[Noreste] & \num{0.844} & \num{0.265} &  & \num{3.19} & \num{0.001} & $[\num{0.326};\ \num{1.363}]$ & \num{0.555} & \num{1.80} \\
Grado universitario (25+) $\times$ region[Medio Oeste] & \num{0.229} & \num{0.187} &  & \num{1.23} & \num{0.220} & $[\num{-0.137};\ \num{0.595}]$ & \num{0.554} & \num{1.80} \\
Grado universitario (25+) $\times$ region[Oeste] & \num{-0.031} & \num{0.192} &  & \num{-0.16} & \num{0.873} & $[\num{-0.408};\ \num{0.346}]$ & \num{0.492} & \num{2.03} \\
Desempleo (16+) $\times$ region[Noreste] & \num{0.903} & \num{0.774} &  & \num{1.17} & \num{0.243} & $[\num{-0.615};\ \num{2.421}]$ & \num{0.791} & \num{1.26} \\
Desempleo (16+) $\times$ region[Medio Oeste] & \num{1.478} & \num{0.273} &  & \num{5.42} & $<\num{0.001}$ & $[\num{0.944};\ \num{2.012}]$ & \num{0.446} & \num{2.24} \\
Desempleo (16+) $\times$ region[Oeste] & \num{-0.429} & \num{0.325} &  & \num{-1.32} & \num{0.187} & $[\num{-1.066};\ \num{0.209}]$ & \num{0.618} & \num{1.62} \\
\end{xltabular}
}


\subsection{Errores típicos clásicos, HC3 y por conglomerados}

Los coeficientes son idénticos con las tres estimaciones de la covarianza; cambian sus errores
típicos (\cref{tab:errores}). Los robustos por conglomerados son, en casi todos los términos,
claramente mayores que los clásicos (la columna «Razón» de la tabla): la inferencia clásica sería
optimista. Con errores típicos clásicos parecerían significativos la natalidad, el bachillerato de
25 años o más, el desempleo en el Sur y la diferencia entre el Medio Oeste y el Sur; con los
robustos no lo son.

{\footnotesize
\begin{xltabular}{\linewidth}{X r r r r r r r}
\caption{Errores típicos del modelo final: clásicos, HC3 y por conglomerados de estado}\label{tab:errores} \\
\toprule
Término & $\hat\beta$ & Clásico & HC3 & Cluster & Razón & $p$ clásico & $p$ cluster \\
\midrule
\endfirsthead
\multicolumn{8}{l}{\emph{(continuación)}} \\
\toprule
Término & $\hat\beta$ & Clásico & HC3 & Cluster & Razón & $p$ clásico & $p$ cluster \\
\midrule
\endhead
\midrule
\multicolumn{8}{r}{\emph{(continúa)}} \\
\endfoot
\bottomrule
\multicolumn{8}{p{0.97\linewidth}}{\footnotesize Razón: E.\,T. por conglomerados entre E.\,T. clásico. Una razón mayor que 1 indica que la inferencia clásica es optimista, por la dependencia dentro de cada estado y la heterocedasticidad.} \\
\endlastfoot
Incidencia de cáncer & \num{0.187} & \num{0.007} & \num{0.012} & \num{0.019} & \num{2.88} & $<\num{0.001}$ & $<\num{0.001}$ \\
Edad mediana & \num{-0.788} & \num{0.095} & \num{0.114} & \num{0.187} & \num{1.97} & $<\num{0.001}$ & $<\num{0.001}$ \\
Tamaño medio del hogar & \num{-12.875} & \num{2.004} & \num{2.758} & \num{3.292} & \num{1.64} & $<\num{0.001}$ & $<\num{0.001}$ \\
Bachillerato (18–24) & \num{0.241} & \num{0.047} & \num{0.061} & \num{0.064} & \num{1.38} & $<\num{0.001}$ & $<\num{0.001}$ \\
Grado universitario (25+) & \num{-1.054} & \num{0.160} & \num{0.183} & \num{0.255} & \num{1.60} & $<\num{0.001}$ & $<\num{0.001}$ \\
Desempleo (16+) & \num{0.366} & \num{0.178} & \num{0.218} & \num{0.284} & \num{1.60} & \num{0.040} & \num{0.204} \\
Solo seguro público & \num{0.892} & \num{0.093} & \num{0.108} & \num{0.214} & \num{2.31} & $<\num{0.001}$ & $<\num{0.001}$ \\
Otra raza & \num{-0.660} & \num{0.118} & \num{0.140} & \num{0.210} & \num{1.78} & $<\num{0.001}$ & \num{0.003} \\
Natalidad & \num{-0.381} & \num{0.187} & \num{0.255} & \num{0.249} & \num{1.33} & \num{0.042} & \num{0.132} \\
region[Noreste] & \num{-15.683} & \num{1.685} & \num{1.294} & \num{2.316} & \num{1.37} & $<\num{0.001}$ & $<\num{0.001}$ \\
region[Medio Oeste] & \num{-4.327} & \num{0.957} & \num{0.950} & \num{2.419} & \num{2.53} & $<\num{0.001}$ & \num{0.080} \\
region[Oeste] & \num{-14.131} & \num{1.221} & \num{1.437} & \num{2.756} & \num{2.26} & $<\num{0.001}$ & $<\num{0.001}$ \\
Bachillerato (25+) & \num{0.354} & \num{0.096} & \num{0.122} & \num{0.192} & \num{1.99} & $<\num{0.001}$ & \num{0.071} \\
Grado universitario (25+) $\times$ region[Noreste] & \num{0.844} & \num{0.265} & \num{0.198} & \num{0.219} & \num{0.83} & \num{0.001} & $<\num{0.001}$ \\
Grado universitario (25+) $\times$ region[Medio Oeste] & \num{0.229} & \num{0.187} & \num{0.188} & \num{0.308} & \num{1.65} & \num{0.220} & \num{0.461} \\
Grado universitario (25+) $\times$ region[Oeste] & \num{-0.031} & \num{0.192} & \num{0.217} & \num{0.257} & \num{1.33} & \num{0.873} & \num{0.905} \\
Desempleo (16+) $\times$ region[Noreste] & \num{0.903} & \num{0.774} & \num{0.546} & \num{0.555} & \num{0.72} & \num{0.243} & \num{0.110} \\
Desempleo (16+) $\times$ region[Medio Oeste] & \num{1.478} & \num{0.273} & \num{0.317} & \num{0.446} & \num{1.64} & $<\num{0.001}$ & \num{0.002} \\
Desempleo (16+) $\times$ region[Oeste] & \num{-0.429} & \num{0.325} & \num{0.462} & \num{0.688} & \num{2.12} & \num{0.187} & \num{0.536} \\
\end{xltabular}
}


\section{Interpretación}\label{sec:interpretacion}

\begin{entregable}{(e): interpretación del modelo}
Cada coeficiente es el cambio esperado en la mortalidad (muertes por \num{100000} habitantes) por
una unidad más de su explicativa, manteniendo constantes las demás. Son asociaciones entre
condados, no efectos causales: el diseño es observacional y ecológico. Para comparar variables
medidas en escalas distintas, la \cref{tab:efectos} expresa cada efecto como el cambio al pasar del
primer al tercer cuartil de la explicativa (rango intercuartílico, RIC).
\end{entregable}

\subsection{La constante y las regiones}

Como las explicativas están centradas, la constante (\res{b-const}) es la mortalidad esperada de un
condado del Sur con todas las explicativas en su media. A igualdad de las demás características, un
condado medio del Noreste tiene \res{b-noreste} muertes por \num{100000} habitantes respecto del
Sur (\res{ic-noreste}), uno del Oeste \res{b-oeste} (\res{ic-oeste}) y uno del Medio Oeste
\res{b-medio-oeste} (\res{ic-medio-oeste}; \res{p-medio-oeste}). La composición socioeconómica
explica una parte importante de las diferencias regionales: la brecha bruta entre el Sur y el Oeste,
\res{brecha-bruta} muertes, se reduce a \res{brecha-ajustada} al ajustar, de modo que el modelo
explica el \res{pct-brecha-explicada}\,\% de ella (\cref{fig:regiones-ajustadas}).

\figura{regiones-ajustadas}{Mortalidad media bruta y ajustada por región. La ajustada es la
predicción del modelo para un condado de esa región con las explicativas en su media.}

\subsection{Efectos principales}

\begin{itemize}
  \item \textbf{Incidencia.} Cada caso diagnosticado más por \num{100000} habitantes se asocia con
  \res{b-incidencia} muertes más \res{ic-incidencia}. Pasar del primer al tercer cuartil
  (\res{iqr-incidencia} casos) supone \res{efiqr-incidencia} muertes más \res{efiqr-incidencia-ic}:
  es, con diferencia, la explicativa más importante ($\text{Beta}=\res{beta-incidencia}$).
  \item \textbf{Cobertura sanitaria sólo pública.} Cada punto porcentual más de residentes cuyo
  único seguro es público (\emph{Medicaid} o \emph{Medicare}) se asocia con \res{b-solopublico}
  muertes más \res{ic-solopublico}; del primer al tercer cuartil, \res{efiqr-solopublico}
  \res{efiqr-solopublico-ic}. Es un indicador de pobreza, edad avanzada y peor acceso a la atención
  especializada.
  \item \textbf{Educación.} El porcentaje de adultos con título universitario tiene el efecto
  protector más claro, pero su magnitud depende de la región (\cref{sec:interacciones-interp}). El
  porcentaje de jóvenes de 18 a 24 años cuyo nivel máximo es el bachillerato se asocia con más
  mortalidad (\res{b-bach1824} por punto, \res{ic-bach1824}), como indicador de menor acceso a
  estudios superiores. El de adultos con bachillerato (\res{b-bach25}, \res{p-bach25}) no es
  significativo, pero se mantiene por ser confusor del efecto del título universitario.
  \item \textbf{Edad mediana.} Cada año más de edad mediana se asocia con \res{b-edad} muertes
  \res{ic-edad}. El signo negativo no significa que envejecer proteja: compara condados con la
  misma incidencia y el mismo perfil socioeconómico. Una explicación compatible, que este diseño no
  permite contrastar, es la migración de jubilados con mejor salud y recursos hacia condados
  envejecidos.
  \item \textbf{Tamaño medio del hogar.} Una persona más por hogar se asocia con \res{b-hogar}
  muertes \res{ic-hogar}; como el rango intercuartílico es sólo de \res{iqr-hogar} personas, el
  efecto entre condados típicos es de \res{efiqr-hogar} \res{efiqr-hogar-ic}.
  \item \textbf{Otra raza.} Cada punto porcentual más de población de otras razas (sobre todo
  hispana, que en el censo puede declararse en esta categoría) se asocia con \res{b-otraraza}
  muertes \res{ic-otraraza}, en línea con la menor mortalidad por cáncer documentada en la población
  hispana.
  \item \textbf{Natalidad} (\res{b-natalidad}, \res{p-natalidad}): no significativa.
\end{itemize}

\begin{table}[htbp]
\centering
\footnotesize
\caption{Magnitud de los efectos: por unidad y al pasar del primer al tercer cuartil}\label{tab:efectos}
\begin{tabularx}{\linewidth}{X l r r r l}
\toprule
Variable & Unidad & Por unidad & IQR & Efecto Q1$\to$Q3 & IC 95\,\% \\
\midrule
Incidencia de cáncer & casos/100 000 hab. & \num{0.187} & \num{66.40} & \num{12.4} & $[\num{9.8};\ \num{15.0}]$ \\
Edad mediana & años & \num{-0.788} & \num{6.10} & \num{-4.8} & $[\num{-7.1};\ \num{-2.5}]$ \\
Tamaño medio del hogar & personas & \num{-12.875} & \num{0.26} & \num{-3.3} & $[\num{-5.1};\ \num{-1.6}]$ \\
Bachillerato (18–24) & \% & \num{0.241} & \num{11.50} & \num{2.8} & $[\num{1.3};\ \num{4.3}]$ \\
Solo seguro público & \% & \num{0.892} & \num{8.25} & \num{7.4} & $[\num{3.8};\ \num{10.9}]$ \\
Otra raza & \% & \num{-0.660} & \num{1.88} & \num{-1.2} & $[\num{-2.0};\ \num{-0.4}]$ \\
Natalidad & \% & \num{-0.381} & \num{1.97} & \num{-0.8} & $[\num{-1.7};\ \num{0.2}]$ \\
Bachillerato (25+) & \% & \num{0.354} & \num{9.25} & \num{3.3} & $[\num{-0.3};\ \num{6.8}]$ \\
\bottomrule
\end{tabularx}
\par\smallskip{\footnotesize\raggedright Muertes por 100\,000 habitantes, manteniendo constantes las demás variables del modelo.\par}
\end{table}


\figura{coeficientes}{Efecto de cada explicativa al pasar del primer al tercer cuartil, con su
intervalo de confianza al 95\,\% (errores típicos por conglomerados). Las variables con interacción
muestran un efecto por región; los marcadores huecos no son significativos.}

\subsection{Interacciones: efectos que dependen de la región}\label{sec:interacciones-interp}

Con interacciones, el coeficiente de una variable en una región es la suma de su efecto principal
(el del Sur) y de la interacción correspondiente; su error típico se obtiene de la combinación
lineal de coeficientes con la matriz de covarianzas robusta. La \cref{tab:pendientes} y la
\cref{fig:interacciones} recogen las pendientes resultantes.

\begin{itemize}
  \item \textbf{Título universitario.} Cada punto porcentual más de adultos con título universitario
  se asocia con \res{pend-grado25-sur} muertes en el Sur (\res{pend-grado25-sur-p}),
  \res{pend-grado25-oeste} en el Oeste (\res{pend-grado25-oeste-p}) y \res{pend-grado25-medio-oeste}
  en el Medio Oeste (\res{pend-grado25-medio-oeste-p}). En el Noreste la pendiente es
  \res{pend-grado25-noreste} y no significativa (\res{pend-grado25-noreste-p}): el gradiente
  educativo es mucho más débil allí.
  \item \textbf{Desempleo.} En el Sur la pendiente es \res{pend-desempleo-sur}
  (\res{pend-desempleo-sur-p}) y en el Oeste \res{pend-desempleo-oeste}
  (\res{pend-desempleo-oeste-p}), ninguna significativa; en cambio, en el Medio Oeste cada punto de
  desempleo se asocia con \res{pend-desempleo-medio-oeste} muertes más
  (\res{pend-desempleo-medio-oeste-p}) y en el Noreste con \res{pend-desempleo-noreste}
  (\res{pend-desempleo-noreste-p}).
\end{itemize}

\begin{table}[htbp]
\centering
\footnotesize
\caption{Efecto de las variables que interaccionan con la región: pendiente específica de cada región}\label{tab:pendientes}
\begin{tabular}{l l r l r r}
\toprule
Variable & Región & Pendiente & IC 95\,\% & $p$ & Efecto Q1$\to$Q3 \\
\midrule
Grado universitario (25+) & Sur & \num{-1.054} & $[\num{-1.566};$\allowbreak\ $\num{-0.542}]$ & $<\num{0.001}$ & \num{-7.1} \\
 & Noreste & \num{-0.209} & $[\num{-0.723};$\allowbreak\ $\num{0.304}]$ & \num{0.416} & \num{-1.4} \\
 & Medio Oeste & \num{-0.825} & $[\num{-1.542};$\allowbreak\ $\num{-0.108}]$ & \num{0.025} & \num{-5.5} \\
 & Oeste & \num{-1.085} & $[\num{-1.600};$\allowbreak\ $\num{-0.569}]$ & $<\num{0.001}$ & \num{-7.3} \\
\addlinespace
Desempleo (16+) & Sur & \num{0.366} & $[\num{-0.205};$\allowbreak\ $\num{0.937}]$ & \num{0.204} & \num{1.5} \\
 & Noreste & \num{1.269} & $[\num{0.221};$\allowbreak\ $\num{2.317}]$ & \num{0.019} & \num{5.3} \\
 & Medio Oeste & \num{1.844} & $[\num{0.940};$\allowbreak\ $\num{2.749}]$ & $<\num{0.001}$ & \num{7.7} \\
 & Oeste & \num{-0.063} & $[\num{-1.417};$\allowbreak\ $\num{1.291}]$ & \num{0.926} & \num{-0.3} \\
\bottomrule
\end{tabular}
\par\smallskip{\footnotesize\raggedright Pendiente en la región $R$: $\hat\beta_{x} + \hat\beta_{x\times R}$, con $\widehat{\mathrm{Var}} = \widehat{\mathrm{Var}}(\hat\beta_x) + \widehat{\mathrm{Var}}(\hat\beta_{x\times R}) + 2\,\widehat{\mathrm{Cov}}$.\par}
\end{table}


\figura{interacciones}{Pendiente del título universitario y del desempleo en cada región, con su
intervalo de confianza al 95\,\%. Marcador hueco: pendiente no significativa.}

\subsection{¿Qué cambia si no se ajusta por la incidencia?}

Incluir la incidencia cambia la pregunta: con ella, los demás coeficientes describen la mortalidad
a igual número de diagnósticos, es decir, sobre todo la supervivencia; sin ella, la asociación
total, que incluye también el riesgo de enfermar. La \cref{tab:incidencia} compara ambos ajustes
sobre la misma muestra. Sin la incidencia el $R^2$ cae de \res{r2-con-incidencia} a
\res{r2-sin-incidencia}, y algunos coeficientes cambian mucho: el del desempleo crece un
\res{cambio-desempleo}\,\%, porque el desempleo se asocia también con mayor incidencia (por ejemplo,
a través del tabaquismo). El del título universitario cambia sólo un \res{cambio-grado25}\,\%: su
asociación con la mortalidad no pasa por la incidencia.

\begin{table}[htbp]
\centering
\footnotesize
\caption{Coeficientes con y sin ajustar por la incidencia (misma muestra y mismos pesos)}\label{tab:incidencia}
\begin{tabularx}{\linewidth}{X r r r}
\toprule
Término & Con incidencia & Sin incidencia & Cambio (\%) \\
\midrule
Edad mediana & \num{-0.788} & \num{-1.240} & \num{-57.3} \\
Tamaño medio del hogar & \num{-12.875} & \num{-21.614} & \num{-67.9} \\
Bachillerato (18–24) & \num{0.241} & \num{0.326} & \num{35.2} \\
Grado universitario (25+) & \num{-1.054} & \num{-0.916} & \num{13.1} \\
Desempleo (16+) & \num{0.366} & \num{0.885} & \num{141.7} \\
Solo seguro público & \num{0.892} & \num{0.884} & \num{-0.9} \\
Otra raza & \num{-0.660} & \num{-1.127} & \num{-70.9} \\
Natalidad & \num{-0.381} & \num{-0.857} & \num{-124.6} \\
Región: Noreste & \num{-15.683} & \num{-9.303} & \num{40.7} \\
Región: Medio Oeste & \num{-4.327} & \num{-4.934} & \num{-14.0} \\
Región: Oeste & \num{-14.131} & \num{-17.237} & \num{-22.0} \\
Bachillerato (25+) & \num{0.354} & \num{0.474} & \num{34.1} \\
\bottomrule
\end{tabularx}
\par\smallskip{\footnotesize\raggedright $R^2$ con incidencia: \num{0.558}; sin ella: \num{0.438}. Un cambio grande indica que la incidencia es confusora (o mediadora) de esa asociación.\par}
\end{table}


\figura{incidencia}{Cambio relativo de cada coeficiente al retirar la incidencia del modelo. Las
líneas discontinuas marcan $\pm 10$\,\%.}
